\section*{Problemas}

\subsection*{Enunciados de los problemas}

% Recordar
\begin{problema}
	\label{prob:43cddbb}
	Enuncie, en símbolos y sin usar ningún dato numérico, la equivalencia exacta entre rechazar $H_0: \theta = \theta_0$ (prueba de dos colas, nivel $\alpha$) y la posición de $\theta_0$ respecto a un intervalo de confianza del $(1-\alpha)100\%$ construido con la misma muestra.
\end{problema}

% Comprender
\begin{problema}
	\label{prob:ae1a813}
	Un ingeniero construye un intervalo de confianza del $95\%$ para la resistencia media de un lote de cables: $\text{IC} = (48.2, 52.6)$ MPa. Usando exclusivamente la equivalencia IC$\leftrightarrow$prueba (sin calcular ningún estadístico de prueba), determine:
	\begin{enumerate}
		\item ¿Se rechazaría $H_0: \mu = 50$ contra $H_a: \mu \neq 50$ al nivel $\alpha = 0.05$?
		\item ¿Se rechazaría $H_0: \mu = 53$ contra $H_a: \mu \neq 53$ al mismo nivel?
	\end{enumerate}
	Justifique cada respuesta únicamente por la posición del valor hipotético respecto al intervalo.
\end{problema}

% Aplicar
\begin{problema}
	\label{prob:1b375b0}
	Una muestra de $n = 25$ mediciones de pH de un proceso químico tiene $\bar{x} = 7.08$ y $s = 0.15$ (varianza poblacional desconocida).
	\begin{enumerate}
		\item Construya el intervalo de confianza del $99\%$ para $\mu$ usando la distribución $t$ ($t_{24,\,0.005} \approx 2.797$).
		\item Use ese intervalo, sin calcular ningún estadístico $T$ adicional, para decidir si se rechaza $H_0: \mu = 7.00$ contra $H_a: \mu \neq 7.00$ al nivel $\alpha = 0.01$.
	\end{enumerate}
\end{problema}

% Analizar
\begin{problema}
	\label{prob:6517bf1}
	Con los mismos datos del problema anterior ($\bar{x} = 7.08$, $s = 0.15$, $n = 25$), un analista realiza una prueba de \textbf{una cola} $H_0: \mu = 7.00$ contra $H_a: \mu > 7.00$ al nivel $\alpha = 0.01$, obtiene $T = 2.667$ y \textbf{rechaza} $H_0$ (ya que $T > t_{24,\,0.01} \approx 2.492$). Sin embargo, en el problema anterior el intervalo de confianza de \textbf{dos colas} del $99\%$ \textbf{no} llevó a rechazar la prueba de dos colas correspondiente. Explique, usando la teoría de esta sección, por qué ambos resultados son consistentes y no contradictorios, identificando específicamente qué valor crítico compara cada procedimiento y por qué son distintos.
\end{problema}

% Evaluar
\begin{problema}
	\label{prob:4f1ea6b}
	Un estudiante concluye que, como $\mu_0 = 50$ cae \textbf{dentro} del intervalo de confianza del $95\%$ del segundo problema de este cuaderno, entonces ``se ha demostrado que la media poblacional real es exactamente 50''. Evalúe esta afirmación: ¿es correcta la conclusión estadística? Relacione su respuesta con la observación de la teoría sobre la diferencia entre ``no rechazar $H_0$'' y ``aceptar $H_0$''.
\end{problema}

% Crear
\begin{problema}
	\label{prob:1e58b99}
	Diseñe un breve ejemplo original (contexto, datos resumidos de una muestra ficticia, e intervalo de confianza calculado a partir de ellos) que ilustre cómo un \textbf{único} intervalo de confianza permite responder simultáneamente tres pruebas de hipótesis distintas, $H_0: \mu = \theta_1$, $H_0: \mu = \theta_2$ y $H_0: \mu = \theta_3$, para tres valores hipotéticos diferentes, sin calcular tres estadísticos de prueba por separado. Presente los tres valores elegidos y la decisión (rechazar/no rechazar) para cada uno, justificada únicamente por su posición respecto al intervalo.
\end{problema}

\subsection*{Soluciones detalladas}

% Recordar
\begin{solproblema}
	\begin{align}
		\text{Rechazar } H_0 \text{ al nivel } \alpha \quad \Longleftrightarrow \quad \theta_0 \notin \text{IC}_{(1-\alpha)100\%}.
	\end{align}
	Es decir, se rechaza $H_0$ si y solo si el valor hipotético $\theta_0$ cae fuera del intervalo de confianza construido con el mismo nivel de confianza $(1-\alpha)100\%$.
\end{solproblema}

% Comprender
\begin{solproblema}
	\begin{enumerate}
		\item $\mu_0 = 50$ está dentro del intervalo $(48.2, 52.6)$, por lo tanto \textbf{no se rechaza} $H_0: \mu = 50$.
		\item $\mu_0 = 53$ está fuera del intervalo $(48.2, 52.6)$ (es mayor que el límite superior $52.6$), por lo tanto \textbf{se rechaza} $H_0: \mu = 53$.
	\end{enumerate}
\end{solproblema}

% Aplicar
\begin{solproblema}
	\begin{enumerate}
		\item Con $\text{SE} = s/\sqrt{n} = 0.15/\sqrt{25} = 0.03$:
		\begin{align}
			\text{IC}_{99\%} = \bar{x} \pm t_{24,\,0.005} \cdot \text{SE} = 7.08 \pm 2.797(0.03) = 7.08 \pm 0.0839 = (6.9961,\ 7.1639).
		\end{align}
		\item Como $\mu_0 = 7.00$ está dentro de $(6.9961,\ 7.1639)$, \textbf{no se rechaza} $H_0: \mu = 7.00$ al nivel $\alpha = 0.01$.
	\end{enumerate}
\end{solproblema}

% Analizar
\begin{solproblema}
	No hay contradicción: cada procedimiento compara el estadístico contra un valor crítico distinto, porque distribuyen el error tipo I de forma diferente.
	\begin{itemize}
		\item La prueba de \textbf{una cola} concentra todo el $\alpha = 0.01$ en una sola cola, por lo que su valor crítico $t_{24,\,0.01} \approx 2.492$ es \textbf{menor} (más fácil de superar).
		\item La prueba de \textbf{dos colas} (equivalente al intervalo de confianza) reparte el $\alpha$ entre las dos colas ($\alpha/2 = 0.005$ en cada una), por lo que su valor crítico $t_{24,\,0.005} \approx 2.797$ es \textbf{mayor} (más difícil de superar).
	\end{itemize}
	El estadístico observado $T = 2.667$ cae \textbf{entre} ambos valores críticos ($2.492 < 2.667 < 2.797$): es lo suficientemente grande para rechazar al nivel de una cola, pero no al nivel de dos colas. Esto es exactamente lo esperado — la equivalencia IC$\leftrightarrow$prueba del teorema de esta sección se enunció explícitamente para pruebas de \textbf{dos colas}; una prueba de una cola no tiene por qué coincidir con la decisión del intervalo de confianza de dos colas, ya que responden preguntas distintas ($\mu \neq \mu_0$ contra $\mu > \mu_0$).
\end{solproblema}

% Evaluar
\begin{solproblema}
	La afirmación es \textbf{incorrecta}. ``No rechazar $H_0$'' únicamente significa que los datos observados son \emph{compatibles} con $\mu_0 = 50$ (no hay evidencia suficiente para descartarlo), pero no constituye una demostración de que $\mu = 50$ sea el valor exacto de la media poblacional. De hecho, cualquier otro valor dentro del intervalo $(48.2, 52.6)$ —por ejemplo $\mu_0 = 49$ o $\mu_0 = 51.5$— tampoco se rechazaría con los mismos datos, por lo que sería lógicamente imposible que todos esos valores fueran simultáneamente ``la media real demostrada''. La estadística inferencial nunca demuestra que un parámetro tome un valor exacto; solo cuantifica qué tan compatibles son los datos con distintos valores hipotéticos.
\end{solproblema}

% Crear
\begin{solproblema}
	\textbf{Ejemplo:} el gerente de control de calidad de una planta embotelladora toma una muestra de $n = 36$ botellas y mide su contenido neto, obteniendo $\bar{x} = 500.4$ ml con $\sigma = 3.0$ ml (conocida, por especificación del equipo de llenado). El intervalo de confianza del $95\%$ para la media de llenado es:
	\begin{align}
		\text{IC}_{95\%} = 500.4 \pm 1.96 \left(\frac{3.0}{\sqrt{36}}\right) = 500.4 \pm 0.98 = (499.42,\ 501.38) \text{ ml}.
	\end{align}
	Con este único intervalo se responden tres pruebas de hipótesis distintas al nivel $\alpha = 0.05$, sin calcular ningún estadístico $Z$ adicional:
	\begin{itemize}
		\item $H_0: \mu = 500$ (el objetivo nominal de llenado): $500 \in (499.42, 501.38)$ $\Rightarrow$ \textbf{no se rechaza}.
		\item $H_0: \mu = 499$ (posible sub-llenado): $499 \notin (499.42, 501.38)$ $\Rightarrow$ \textbf{se rechaza}.
		\item $H_0: \mu = 502$ (posible sobre-llenado): $502 \notin (499.42, 501.38)$ $\Rightarrow$ \textbf{se rechaza}.
	\end{itemize}
	Las tres decisiones se obtuvieron únicamente comparando cada valor hipotético contra los mismos dos límites del intervalo.
\end{solproblema}
